% Part: proof-theory
% Chapter: sequent-calculus
% Section: rules-proofs
%READERNOTE{SOL6-A336: सिद्धतेच्या व्याख्येतील विगमनाची संज्ञा सुसंगत केली आणि उघडे अवतरणचिन्ह बंद केले.}

\documentclass[../../../include/open-logic-section]{subfiles}

\begin{document}

\olfileid{pt}{seq}{rp}

\olsection{नियम आणि \usetoken{P}{proof}}

आधी काही व्याख्या देऊ आणि काही संज्ञांचा अर्थ निश्चित करू.

उदाहरणार्थ, $!A$ आणि~$!B$ असलेल्या क्रमवर्तींपासून
$!A \land !B$ असलेल्या क्रमवर्ती !!{proving} करू देणारे \Log{G3c} मधील
दोन नियम पाहा:
\[
\Axiom$ !A, !B, \Gamma \fCenter \Delta$
\RightLabel{\LeftR{\land}}
\UnaryInf$ !A \land !B, \Gamma \fCenter \Delta$
\DisplayProof
\qquad
\Axiom$\Gamma \fCenter \Delta, !A$
\Axiom$ \Gamma \fCenter \Delta, !B$
\RightLabel{\RightR{\land}}
\BinaryInf$ \Gamma \fCenter \Delta, !A \land !B $
\DisplayProof
\]
रेषेच्या वरच्या क्रमवर्तींना \emph{आधारविधाने}, तर खालच्या
क्रमवर्तीला \emph{निष्कर्ष} म्हणतात. प्रत्येक नियमातील
$\Gamma$ आणि $\Delta$ यांतील !!{formula}s ना \emph{संदर्भ}
किंवा \emph{बाजूची !!{formula}s} म्हणतात. निष्कर्षातील
संदर्भाबाहेरचे !!{formula} हे \emph{मुख्य} !!{formula} असते;
आणि आधारविधानांतील संदर्भाबाहेरच्या !!{formula}s ना
\emph{सक्रिय} (किंवा \emph{सहायक}) !!{formula}s म्हणतात.

परिमाणकांचे नियम थोडे अधिक गुंतागुंतीचे आहेत. उदाहरणार्थ,
$\lforall$ साठीचे \Log{G1c} मधील नियम असे आहेत:
\[\Axiom$ !A(t), \Gamma \fCenter \Delta$
\RightLabel{\LeftR{\lforall}}
\UnaryInf$ \lforall[x][!A(x)],\Gamma \fCenter \Delta$
\DisplayProof
\qquad
\Axiom$ \Gamma \fCenter \Delta, !A(a) $
\RightLabel{\RightR{\lforall}}
\UnaryInf$ \Gamma \fCenter \Delta, \lforall[x][!A(x)]$
\DisplayProof
\]
\LeftR{\lforall} नियमातील सक्रिय !!{formula} हे~$!A(t)$ आहे;
येथे $t$ हे बंद पद, म्हणजे चर नसलेले पद, आहे. एखादे
निगमन योग्य असते—म्हणजे ते \LeftR{\lforall} या
रूपरेषात्मक नियमाशी जुळते—जर एखादे !!{formula}~$!A$,
चर~$x$ आणि बंद पद~$t$ असे असतील की निष्कर्षातील मुख्य
!!{formula} $\lforall[x][!A]$ असेल आणि आधारविधानातील
सक्रिय !!{formula} हे~$\Subst{!A}{t}{x}$ असेल.
\RightR{\lforall} नियमही याचप्रमाणे कार्य करतो; मात्र
कोणत्याही $t$ पदाऐवजी आधारविधानातील सक्रिय !!{formula}
$\Subst{!A}{a}{x}$ असे असले पाहिजे, जेथे $a$ हा
!!a{constant} आहे. या अचराला \RightR{\lforall}
निगमनाचा \emph{आयगेन चर} म्हणतात. खालच्या क्रमवर्तीमध्ये
$a$ येत नसेल—म्हणजेच $\Gamma$, $\Delta$ किंवा~$!A$
यांत $a$ नसेल—तरच हे निगमन योग्य असते. ही
\emph{आयगेन चराची अट} होय.

\Log{G1c} आणि \Log{G3c} यांसारख्या क्रमवर्ती पद्धतींत अशा
रूपरेषात्मक रीतीने दिलेल्या नियमांचा संग्रह आणि स्वयंसिद्धके
म्हणून ग्राह्य अशा, त्याच रीतीने दिलेल्या, काही क्रमवर्ती असतात.
अशा पद्धतीतील !!^{proof}s म्हणजे स्वयंसिद्धकांपासून निगमन नियमांनी
विगमनाने तयार होणारे क्रमवर्तींचे सान्त वृक्ष. प्रत्येक पान
स्वयंसिद्धक असते; आणि इतर प्रत्येक क्रमवर्ती तिच्या लगेच
वरच्या क्रमवर्तींपासून पद्धतीतील एका निगमन नियमाने मिळते.

\begin{defn}\ollabel{defn:proof}
क्रमवर्ती पद्धतीतील \emph{!!^{proof}s} म्हणजे खालीलप्रमाणे
विगमनाने बनवलेले क्रमवर्तींचे सान्त वृक्ष:
\begin{enumerate}
  \item\ollabel{defn:prf:axiom} प्रत्येक स्वयंसिद्धक ही !!a{proof} असते.
  \item\ollabel{defn:prf:one-prem} $\pi_1$ ही $S_1$ या क्रमवर्तीवर
  संपणारी !!a{proof} असेल आणि $S_1$ पासून एका आधारविधानाच्या
  $R$ नियमाने $S$ ही क्रमवर्ती मिळत असेल, तर
  \[
  \AxiomC{}
  \RightLabel{$\pi_1$}
  \DeduceC{$S_1$}
  \RightLabel{$R$}
  \UnaryInfC{$S$}
  \DisplayProof
  \]
  ही !!a{proof} असते.
  \item\ollabel{defn:prf:two-prem} $\pi_1$ आणि $\pi_2$ या
  अनुक्रमे $S_1$ आणि~$S_2$ या क्रमवर्तींवर संपणाऱ्या !!{proof}s
  असतील, आणि $S_1$ व $S_2$ यांच्यापासून दोन आधारविधानांच्या $R$ नियमाने
  $S$ मिळत असेल, तर
  \[
  \AxiomC{}
  \RightLabel{$\pi_1$}
  \DeduceC{$S_1$}
  \AxiomC{}
  \RightLabel{$\pi_2$}
  \DeduceC{$S_2$}
  \RightLabel{$R$}
  \BinaryInfC{$S$}
  \DisplayProof
  \]
  ही !!a{proof} असते.
\end{enumerate}
\end{defn}

\Olref{defn:proof} मध्ये !!{proof}s ची व्याख्या क्रमवर्तींचे वृक्ष म्हणून
दिली आहे. हे वृक्ष ``वरच्या दिशेने'' वाढतात असे आपण मानतो.
सर्वांत वरची (पानांवरील) स्थाने स्वयंसिद्धके असतात; त्यांना
\emph{आरंभीच्या क्रमवर्ती} म्हणतात. सर्वांत खालची क्रमवर्ती
\emph{अंतिम क्रमवर्ती} असते. $\pi$ ही $S$ अंतिम क्रमवर्ती
असलेली !!a{proof} असेल तर $\pi \Proves S$ असेही लिहितो.
$S$ ची !!a{proof} असेल तर $\Proves S$ असे लिहितो; काही वेळा
ज्या पद्धतीत सिद्धता आहे तीही $\Proves$ च्या डावीकडे दाखवतो,
उदा., $\Log{G1c} \Proves !A, !B \Sequent !A \land
!B$. आरंभीच्या क्रमवर्ती वगळता !!a{proof} मधील प्रत्येक
क्रमवर्ती एखाद्या $R$ नियमाने झालेल्या निगमनाचा
\emph{निष्कर्ष} असते. !!a{proof} मधील एखाद्या क्रमवर्तीच्या
लगेच वर असलेल्या एक किंवा दोन क्रमवर्तींना (वृक्षातील तिच्या
पालक-स्थानांना) त्या निगमनाची \emph{आधारविधाने} म्हणतात.

एखादी !!a{proof} विगमनाने बनवताना वापरलेल्या !!{proof}s ना तिच्या
\emph{उप-!!{proof}s} म्हणतात. एखाद्या निगमनाच्या आधारविधानांवर
संपणाऱ्या उप-!!{proof}s ना त्या निगमनाच्या \emph{लगेचच्या
उप-!!{proof}s} म्हणतात.

अनेकदा !!{proof}s ची गुंतागुंत एखाद्या नैसर्गिक संख्येने
मोजावी लागते. !!a{proof} मधील क्रमवर्तींची संख्या मोजता येईल;
पण अनेकदा अधिक उपयोगी मोजमाप म्हणजे !!a{proof} ची !!{height}:
अंतिम क्रमवर्तीपासून आरंभीच्या क्रमवर्तीपर्यंतच्या सर्वाधिक
निगमन-पायऱ्यांची संख्या. तिची विगमनात्मक व्याख्याही देता येते.

\begin{defn}
!!a{proof}~$\pi$ ची \emph{!!{height}}~$\pheight{\pi}$ पुढीलप्रमाणे
विगमनाने ठरते.
\begin{enumerate}
  \item $\pi$ मध्ये एकच स्वयंसिद्धक असेल, तर $\pheight{\pi} = 0$.
  \item $\pi$ चा शेवट एका आधारविधानाच्या निगमनाने होत असेल आणि
  तिची लगेचची उप-!!{proof}~$\pi_1$ असेल, तर $\pheight{\pi} = \pheight{\pi_1} + 1$.
  \item $\pi$ चा शेवट दोन आधारविधानांच्या निगमनाने होत असेल आणि
  तिच्या लगेचच्या उप-!!{proof}s~$\pi_1$ आणि~$\pi_2$ असतील, तर $\pheight{\pi} =
  \max(\pheight{\pi_1}, \pheight{\pi_2}) + 1$.
\end{enumerate}
$S$ या क्रमवर्तीची $\le n$ उंचीची !!a{proof} असेल, तर $\Proves[n] S$ असे लिहितो.
\end{defn}

\begin{ex}
\Log{G3c} मध्ये $!C, \Gamma \Sequent \Delta, !C$ या रूपाची
प्रत्येक क्रमवर्ती स्वयंसिद्धक असते; येथे $!C$ अण्वीय असले पाहिजे.
$!D$ आणि $!E$ ही अण्वीय वाक्ये असोत. मग $!D, !E \Sequent !D$
आणि $!D, !E \Sequent !E$ ही $!C, \Gamma \Sequent \Delta, !C$
या स्वयंसिद्धकाची उदाहरणे आहेत. उदाहरणार्थ, $!C \ident !E$,
$\Gamma = \{!D\}$ आणि $\Delta = \emptyset$ घेतल्यास
$!C, \Gamma \Sequent \Delta, !C$ पासून नेमकी $!D, !E
\Sequent !E$ ही क्रमवर्ती मिळते. (लक्षात घ्या, आपण \emph{बहुसंच} वापरतो;
म्हणून $!D, !E \Sequent !E$ आणि $!E, !D \Sequent !E$ या एकाच
क्रमवर्ती आहेत.)

$!D, !E \Sequent !D$ ही \Log{G3c} च्या स्वयंसिद्धकाचे उदाहरण
असल्यामुळे ती स्वतःच \Log{G3c} मधील !!a{proof} असते. त्याचप्रमाणे
$!D, !E \Sequent !E$ हीही \olref{defn:proof}\olref{defn:prf:axiom}
नुसार तशीच असते. त्यांना अनुक्रमे $\pi_1$ आणि~$\pi_2$ म्हणू.
\olref{defn:prf:two-prem} नुसार, या दोन !!{proof}s घेऊन
दोन आधारविधानांच्या $\RightR{\land}$ नियमाने नवी सिद्धता
($\pi_3$) तयार करता येते:
\[
\Axiom$!D, !E \fCenter !D$
\Axiom$!D, !E \fCenter !E$
\RightLabel{\RightR{\land}}
\BinaryInf$!D, !E  \fCenter !D \land !E $
\DisplayProof
\]
$\pi_3$ पासून पुढे !!{proof}~$\pi_4$ मिळते:
\[
\Axiom$!D, !E \fCenter !D$
\Axiom$!D, !E \fCenter !E$
\RightLabel{\RightR{\land}}
\insertBetweenHyps{\hspace{5ex}}
\BinaryInf$!D, !E  \fCenter !D \land !E $
\LeftSubproofLabel{$\pi_3$}
\RightLabel{\LeftR{\land}}
\UnaryInf$!D \land !E  \fCenter !D \land !E$
\RightSubproofLabel{$\pi_4$}
\DisplayProof
\]
येथे एका आधारविधानाचा $\LeftR{\land}$ नियम वापरला आहे
(\olref{defn:proof}\olref{defn:prf:one-prem} नुसार). $\pi_4$ मधील
शेवटच्या निगमनाची लगेचची उप-!!{proof} ही~$\pi_3$ आहे.
$\RightR{\land}$ निगमनाच्या लगेचच्या उप-!!{proof}s $\pi_1$
आणि~$\pi_2$ आहेत. $\pi_4$ च्या उप-!!{proof}s मध्ये $\pi_1$,
$\pi_2$, $\pi_3$ आणि स्वतः $\pi_4$ यांचा समावेश होतो.
$\pheight{\pi_1} = \pheight{\pi_2} = 0$,
$\pheight{\pi_3} = 1$ आणि $\pheight{\pi_4} = 2$.
म्हणून कोणत्याही अण्वीय $!D$ आणि~$!E$ साठी
$\Log{G3c} \Proves[2] !D \land !E  \fCenter !D
\land !E$ हे दाखवले आहे.
\end{ex}

\end{document}
